\documentclass[../../../uem_mat_tok_q.tex]{subfiles}

\begin{document}
    整式 $f(x)=(x-1)^2(x-2)$ を考える．

    \begin{enumerate}
        \item $g(x)$ を実数を係数とする整式とし，$g(x)$ を $f(x)$ で割った余りを $r(x)$ とおく．%
            $g(x)^7$ を $f(x)$ で割った余りと $r(x)^7$ を $f(x)$ で割った余りが等しいことを示せ．
        \item $a$，$b$ を実数とし，$h(x)=x^2+ax+b$ とおく．%
            $h(x)^7$ を $f(x)$ で割った余りを $h_1(x)$ とおき，%
            $h_1(x)^7$ を $f(x)$ で割った余りを $h_2(x)$ とおく．%
            $h_2(x)$ が $h(x)$ に等しくなるような $a$，$b$ の組をすべて求めよ．
    \end{enumerate}
\end{document}

